\chapter{Proprietary Market Makers and High-Frequency Firms}\label{ch:in:proprietary-market-makers-and-high-frequency-firms}

One electronic market maker files a Form 10-K every February: its revenue, its costs, its headcount and its offices are
public. Firms of comparable or larger size file nothing of the kind. What the public can learn about them comes through
side doors: the one page of their United States broker-dealer's annual report that the law makes public, the accounts
of their British subsidiaries at Companies House, their own web pages, the talks and code their engineers publish, and
what regulators find when they study the market as a whole. This chapter reads those doors for a handful of
proprietary market makers and high-frequency firms, and builds a profile of each that says what is known, as of when,
and from where --- and, as plainly, what is not.

\section{What the public record holds for a private trading firm}

A private trading firm has no shareholders to report to, but it is not invisible. Four kinds of record are public.
\begin{itemize}
\item \textbf{The broker-dealer's statement.} A US broker-dealer files an annual report with the SEC; the statement of
financial condition, its notes and the auditor's report are filed as a separate document that is ``not confidential'',
while the income statement and the periodic reports may be kept confidential. The public part gives the entity's
balance sheet and a paragraph describing its business.
\item \textbf{Subsidiaries' accounts.} A company registered in the United Kingdom files audited accounts every year,
free to download, with a strategic report describing its activity (chapter~11).
\item \textbf{The firm's own publications.} Home pages, careers pages, engineering blogs, open-source code, conference
talks.
\item \textbf{Regulators' and researchers' studies} of the market, which describe the industry without naming firms.
\end{itemize}

\begin{definition}[Firm profile]\label{def:in:proprietary-market-makers-and-high-frequency-firms:profile}
A \emph{firm profile}\index{firm profile}, in this book, is a set of facts about an employer, each a range with a unit,
the date it describes, its scope (the whole firm or one legal entity of it) and the public source it comes from. A field
with no admissible fact is left empty and says so.
\end{definition}

The scope matters most. A subsidiary's accounts describe the subsidiary, not the group: a British entity that trades on
behalf of its group under a transfer-pricing agreement (Book~16, chapter~15) earns what the agreement assigns it.

\begin{example}[One group's European subsidiary]\label{ex:in:proprietary-market-makers-and-high-frequency-firms:jse}
Jane Street Europe Limited, a London-based subsidiary of Jane Street Group, LLC, states in its 2025 strategic report
that its group of subsidiaries ``primarily carries out principal trading activities'' and is compensated for trading
performed on behalf of the group under a trading services agreement and the group's transfer-pricing arrangements. It
reports revenues of \$643.0 million for 2025 against \$995.8 million for 2024, a fall of 35.4\%, after-tax profits of
\$437.4 million (68.0\% of revenue) against \$654.9 million, and total equity of \$3\,058.9 million, up 15.1\%. These
are the numbers of one entity of the group under the group's own allocation rules; they are not the firm's revenue, and
they cannot be scaled up to it.
\end{example}

\begin{example}[One broker-dealer's balance sheet]\label{ex:in:proprietary-market-makers-and-high-frequency-firms:cdrg}
The public statement of financial condition of Citadel Securities LLC at 31 December 2023 describes a company that
``primarily engages in market making and liquidity provision in U.S.\ options, equities, government securities, and
foreign exchange products'', with total assets of \$52\,344 million and member's capital of \$4\,693 million: assets of
11.2 times capital. It says nothing of revenue, profit or staff.
\end{example}

\begin{method}[Building a profile from public records]\label{met:in:proprietary-market-makers-and-high-frequency-firms:build}
\begin{enumerate}
\item List the legal entities (chapter~1, method 1.1) and, for each, the public filings that exist.
\item For each field of the profile, take the most primary source: a filing before the firm's page, the firm's page
before the press.
\item Record every fact as a range with a unit, a date and a scope; a figure the source rounds is the range it rounds
(``300 employees'' is 250--350).
\item Keep the entity's figures apart from the firm's: never add, scale or extrapolate across entities.
\item Leave a field empty, and say so, when no source gives it.
\end{enumerate}
\end{method}

\section{Origins: from the floors to the screens}

The firms of this chapter that state a founding year give one in the 1990s or early 2000s, the years in which
trading moved from exchange floors to screens.

\begin{definition}[Open outcry, floor trader]\label{def:in:proprietary-market-makers-and-high-frequency-firms:floor}
\emph{Open outcry}\index{open outcry} is trading on an exchange floor or in a pit, where members make bids and offers
by voice and hand signals and trade face to face. A \emph{floor trader}\index{floor trader} is an exchange member who
traded there in person, for its own account (a ``local'') or for others; the US industry classification still lists
``securities floor traders'' and ``commodity contracts floor traders'' as principals dealing with investors.
\end{definition}

The floors are almost gone. The largest US derivatives exchange group reports that 93\% of its 2025 volume was
electronic, with \omterm{def:in:proprietary-market-makers-and-high-frequency-firms:floor}{open outcry} kept in Chicago only for options on SOFR futures
(\cref{fig:in:proprietary-market-makers-and-high-frequency-firms:cme}). What the floors left to the firms that
replaced them is a way of working: quote both sides, hold the position briefly, hedge it, and measure the day's result
at the close. \Cref{fig:in:proprietary-market-makers-and-high-frequency-firms:timeline} places the founding dates the
firms themselves give.

\begin{omfigure}
\begin{tikzpicture}
\begin{axis}[width=10.5cm,height=5cm,ybar stacked,bar width=26pt,xmin=-0.6,xmax=1.6,ymin=0,ymax=30,
  xtick={0,1},xticklabels={2024,2025},ylabel={\small contracts a day (millions)},
  tick label style={font=\footnotesize},grid=major,grid style={gray!20},table/col sep=comma,
  legend style={legend cell align=left,font=\scriptsize,at={(0.5,-0.25)},anchor=north,/tikz/every even column/.append style={column sep=8pt}},legend columns=3]
\addplot[fill=omDef!55,draw=omDef] table[x=k,y=globex]{figdata/industry/02-proprietary-market-makers-and-high-frequency-firms/cme.csv};
\addplot[fill=omThm!55,draw=omThm] table[x=k,y=open_outcry]{figdata/industry/02-proprietary-market-makers-and-high-frequency-firms/cme.csv};
\addplot[fill=gray!40,draw=gray] table[x=k,y=privately_negotiated]{figdata/industry/02-proprietary-market-makers-and-high-frequency-firms/cme.csv};
\legend{electronic (Globex),open outcry,privately negotiated}
\end{axis}
\end{tikzpicture}

\omcaption{Where the largest US derivatives exchange group's contracts traded, 2024 and 2025: 24.5 and 26.2 million
contracts a day on its electronic platform, 1.0 and 0.9 million by \omterm{def:in:proprietary-market-makers-and-high-frequency-firms:floor}{open outcry} (options on SOFR futures only), 1.0
million privately negotiated. Data: CME Group Form 10-K for 2025, through \texttt{in\_profiles.cme}.}\label{fig:in:proprietary-market-makers-and-high-frequency-firms:cme}
\end{omfigure}

\begin{omfigure}
\begin{tikzpicture}[font=\footnotesize,>=stealth]
\draw[->,thick,gray] (0,0) -- (12.2,0);
\foreach \x/\y in {0/1990,2/1994,4/1998,6/2002,8/2006,10/2010}
  {\draw[gray] (\x,0.08) -- (\x,-0.08); \node[below,gray] at (\x,-0.1) {\y};}
\foreach \x/\h/\t in {1/2.8/{DRW (1992)},4/2.1/{Tower Research (1998)},5/1.4/{Jane Street (2000)},6/0.7/{Hudson River Trading (2002)}}
  {\draw[omDef,line width=0.9pt] (\x,0) -- (\x,\h); \fill[omDef] (\x,0) circle (1.6pt); \node[anchor=south west,inner sep=2pt] at (\x,\h) {\t};}
\end{tikzpicture}

\omcaption{Founding years that four proprietary trading firms give on their own pages. The period is the one in which
equity and futures exchanges moved from floors to screens. Each date is the firm's own; firms that state none are
left out. Sources: the firms' pages (ledger F2, F5, F7,
F9).}\label{fig:in:proprietary-market-makers-and-high-frequency-firms:timeline}
\end{omfigure}

\begin{remark}[What a founding year does not say]\label{rem:in:proprietary-market-makers-and-high-frequency-firms:founding}
A founding year dates a legal entity or a partnership, not the firm's current business. One firm states that it made
its first bitcoin trade in 2010 and launched a crypto trading business in 2014; another that it launched a client
market-making business in 2022. Firms of this kind change products faster than names.
\end{remark}

\section{Profiles: products, headcount and offices}

The dated box gathers what the firms and their filings state. Each cell is a range from one source, dated; an empty
cell means that no primary source gave the field on the date of access.

\begin{dated}{2026-09}{Eight profiles from public sources}\label{dat:in:proprietary-market-makers-and-high-frequency-firms:profiles}
\small
\begin{tabular}{@{}lllll@{}}
\toprule
firm & founded & employees & offices & firm revenue\\
\midrule
DRW & 1992 & over 2\,000 & -- & --\\
Hudson River Trading & 2002 & -- & 14 cities & --\\
Jane Street & 2000 & over 3\,000 & 5 & --\\
Jump Trading & -- & -- & 13 cities & --\\
Tower Research & 1998 & -- & -- & --\\
Virtu Financial & -- & about 1\,027 & 13 countries & \$3\,632 million (2025)\\
XTX Markets & -- & 250--350 & -- & --\\
Citadel Securities & -- & -- & -- & --\\
\bottomrule
\end{tabular}

\smallskip\noindent Firms' own statements of what they do: ``a quantitative trading firm and liquidity provider''
(Jane Street); ``a multi-asset class quantitative trading firm that provides liquidity on global markets and directly
to our clients'' (Hudson River Trading); ``a home for \dots\ quantitative trading teams, powered by a high-performance
technology platform'' (Tower Research); ``a diversified trading firm'' (DRW); ``a global trading firm'' in ``every
asset class'' (Jump Trading); an algorithmic trading firm producing ``price forecasts for over 50\,000 financial
instruments'', with ``\$250bn daily traded volume'' (XTX Markets). Virtu Financial reports a market-making and an
execution-services segment. For Citadel Securities only its US broker-dealer's balance sheet is public
(\cref{ex:in:proprietary-market-makers-and-high-frequency-firms:cdrg}).
\end{dated}

Three things stand out. Only the listed firm has a revenue figure; for the others, the revenue numbers that circulate
come from press accounts of documents shown to lenders, which are not public, and this book does not use them. The
headcounts that are public are lower bounds or rounded figures: ``over 3\,000'' is a range with no upper end. And
offices are counted in different units, cities, countries or offices, which cannot be compared without the list.

\begin{remark}[Reading headcounts]\label{rem:in:proprietary-market-makers-and-high-frequency-firms:heads}
A firm that states ``more than 3\,000 employees'' in 2026 and a filing that states ``approximately 1\,027 employees''
describe two firms whose revenue per head is unknown for one and \$3.5 million for the other. Revenue per head is chapter~12's subject; it can be computed only where both numbers exist for
the same scope and year.
\end{remark}

\section{Culture and technology as the firms describe them}

The firms publish a great deal about how they work, because they recruit on it. Their statements are the firms'
own, not independent evidence, and are read as statements of what the firm wants a candidate to know.
\begin{itemize}
\item \textbf{Languages and code.} One firm states that it uses OCaml, ``a statically typed functional language, as our
primary development platform''. Another's crypto arm publishes Firedancer, ``a new Solana validator client written
entirely in C''. The choice of language is also a statement about hiring: it tells a candidate what will be learned on
the job (Book~13, chapter~10).
\item \textbf{Research as scale.} One firm describes ``an unrivalled level of computational resources \dots\ with a
growing research cluster'' and machine-learning forecasts for tens of thousands of instruments; another lists ``800+
technologists'' among over 2\,000 staff.
\item \textbf{Organisation.} One firm describes itself as ``a home for \dots\ quantitative trading teams'' on a shared
technology platform: a structure closer to chapter~5's platforms than to a single-book market maker.
\item \textbf{Clients.} Firms that began trading only against the market now make prices to clients directly: one dates
its single-dealer platform to 2018 and its client market-making to 2022.
\end{itemize}

\section{Size and concentration}

Regulators' studies of the market describe the industry's concentration without naming firms. The UK regulator's study
of latency-arbitrage races on the London Stock Exchange found ``clear concentration of winners, with the top 3 firms
winning 54\% of races, and the top 6 firms winning 82\% of races''; the same six firms lost 85\% of races, so the
business of racing is a contest among a handful of firms on both sides. Races took about 22\% of the daily volume of
the largest UK shares, and the study put the sums at stake at about \$5 billion a year across global equity markets
(Book~11, chapter~9, describes the races).

For an employee the concentration has two readings. The largest firms in the fastest businesses are few, so a
specialist's market is small and its employers know one another. And the firms that win races are the ones that also
lose them: what they compete on is speed and pricing at the margin, which is why the engineering chapters of Part~III
describe their jobs as they do.

\section{Tutorial: eight profiles and their gaps}\label{tut:in:proprietary-market-makers-and-high-frequency-firms:tut}

\begin{tutorial}
\textbf{Goal.} Build the chapter's profiles from sourced facts and measure what public sources leave out.
\textbf{End state:} the dated box's table and \cref{fig:in:proprietary-market-makers-and-high-frequency-firms:cov}.
\begin{tutsteps}
\item \textbf{The facts.} \texttt{data/industry/profiles/facts.csv} holds 24 facts, one per row: firm, field, low and
high, unit, date, scope (firm or entity), ledger row, and a text value for descriptive fields.
\item \textbf{Validate.} \texttt{in\_profiles.check()} runs \texttt{firm.profiles.problems} against the ledger rows of
this chapter (\cref{lst:in:proprietary-market-makers-and-high-frequency-firms:problems}); it returns an empty list.
\omcode{code/firm/profiles/firm_profiles.py}{53}{70}{A fact enters a profile only with a ledger row, a date, an ordered
range and a value.}\label{lst:in:proprietary-market-makers-and-high-frequency-firms:problems}
\item \textbf{Render.} \texttt{firm.profiles.table(facts, fields)} prints one row per firm; entity-scope facts (the
subsidiary's revenue, the broker-dealer's assets) are kept out of the firm columns.
\item \textbf{Measure the gap.} \texttt{in\_profiles.gap()} counts, for five standard fields (founded, employees,
offices, firm revenue, products), how many each firm's public sources fill.
\end{tutsteps}
\end{tutorial}

\begin{omfigure}
\begin{tikzpicture}
\begin{axis}[width=10.5cm,height=5.6cm,xbar stacked,bar width=8pt,xmin=0,xmax=5,xtick={0,1,2,3,4,5},xlabel={\small fields filled from public sources (of 5)},
  ytick={0,1,2,3,4,5,6,7},yticklabels from table={figdata/industry/02-proprietary-market-makers-and-high-frequency-firms/coverage.csv}{firm},
  y dir=reverse,tick label style={font=\footnotesize},grid=major,grid style={gray!20},table/col sep=comma,
  legend style={legend cell align=left,font=\scriptsize,at={(0.5,-0.25)},anchor=north,/tikz/every even column/.append style={column sep=8pt}},legend columns=2]
\addplot[fill=omDef!55,draw=omDef] table[y=k,x=listed]{figdata/industry/02-proprietary-market-makers-and-high-frequency-firms/coverage.csv};
\addplot[fill=omThm!55,draw=omThm] table[y=k,x=private]{figdata/industry/02-proprietary-market-makers-and-high-frequency-firms/coverage.csv};
\legend{listed firm,private firm}
\end{axis}
\end{tikzpicture}

\omcaption{The disclosure gap: of five standard profile fields (founded, employees, offices, firm revenue, products),
how many each firm's public sources fill at the firm level on 29 September 2026. The listed firm fills four; the private
firms fill 0 to 4, median 2, and none fills firm revenue. Data: \texttt{data/industry/profiles/facts.csv}, through
\texttt{in\_profiles.gap}.}\label{fig:in:proprietary-market-makers-and-high-frequency-firms:cov}
\end{omfigure}

The listed firm is missing only a founding year that its annual report does not state. A private firm with an active
careers site fills as many fields; the others fill two or fewer, and the broker-dealer whose balance sheet is public
fills none at the firm level. Revenue is the field that separates them: no private firm in the sample has a primary
source for it.

\textbf{What to change next.} Add the UK subsidiaries' average number of employees as entity-scope facts and see why
they cannot fill the firm's headcount (exercise~7); tighten the staleness limit to 60 days and see which facts would
have to be re-verified before a reprint.

\section{Build: the profile store}\label{bld:in:proprietary-market-makers-and-high-frequency-firms:profiles}

\begin{build}
\bfield{Purpose} One way to hold every firm fact the book prints, so that chapters 3--10 and 12 check each fact against
the ledger before it reaches the page.
\bfield{Interface} \texttt{firm.profiles}: \texttt{Fact(firm, field, low, high, unit, as\_of, scope, ledger, text)};
\texttt{load\_facts(path)}; \texttt{problems(facts, ledger\_ids, today, max\_age\_days)}; \texttt{widen(value,
rounding)}; \texttt{coverage(facts, fields, scope)}; \texttt{gap(facts, fields, listed)}; \texttt{table(facts, fields)}.
\bfield{Rules} No fact without a ledger row and a date; a point value only where the source gives one, a rounded figure
as the range it rounds; entity facts never fill a firm field; an open bound stays open.
\bfield{Acceptance tests} \texttt{code/firm/profiles/tests/}: an admissible fact passes; each rule fires on a fact that
breaks it; \texttt{widen} and \texttt{coverage} on constructed facts; table cells print years without separators.
\bfield{Stretch} Read facts straight from the source ledger's rows; flag facts older than the book's 60-day rule for
dated boxes.
\end{build}

\begin{omsources}
\item Virtu Financial, Form 10-K for 2025; CME Group, Form 10-K for 2025.
\item Citadel Securities LLC, Form X-17A-5 Part III, statement of financial condition at 31 December 2023 (public
document under Rule 17a-5(e)(3)).
\item Companies House: annual reports for 2025 of Jane Street Europe Limited (05903707), Hudson River Trading Europe
Ltd.\ (06796079) and XTX Markets Limited (09415174).
\item The firms' own pages (ledger F2--F15), accessed 29 September 2026.
\item M.~Aquilina, E.~Budish and P.~O'Neill, \emph{Quantifying the High-Frequency Trading ``Arms Race''}, FCA Occasional
Paper 50, 2020.
\end{omsources}

\section{Exercises}

\begin{exercise}[$\star$]\label{exo:in:proprietary-market-makers-and-high-frequency-firms:1}
From the CME figures, what share of the exchange group's 2025 volume was traded by \omterm{def:in:proprietary-market-makers-and-high-frequency-firms:floor}{open outcry}, and by how much did
open-outcry volume fall from 2024?
\end{exercise}

\begin{exercise}[$\star$]\label{exo:in:proprietary-market-makers-and-high-frequency-firms:2}
Write the range each statement implies: ``300 employees'', ``more than 3\,000 employees'', ``approximately 1\,027
employees''. Which has no upper bound?
\end{exercise}

\begin{exercise}[$\star$]\label{exo:in:proprietary-market-makers-and-high-frequency-firms:3}
From \cref{ex:in:proprietary-market-makers-and-high-frequency-firms:cdrg}, compute the broker-dealer's member's capital
and its ratio of assets to capital.
\end{exercise}

\begin{exercise}[$\star\star$]\label{exo:in:proprietary-market-makers-and-high-frequency-firms:4}
From \cref{ex:in:proprietary-market-makers-and-high-frequency-firms:jse}, compute the subsidiary group's after-tax
margin in 2024 and 2025 and its return on average equity in 2025. Why can none of these be read as the firm's?
\end{exercise}

\begin{exercise}[$\star\star$]\label{exo:in:proprietary-market-makers-and-high-frequency-firms:5}
In the race study, the top six firms win 82\% of races and lose 85\%. What share of races do they win and lose among
themselves at least, and what does that say about who they compete with?
\end{exercise}

\begin{exercise}[$\star\star$]\label{exo:in:proprietary-market-makers-and-high-frequency-firms:6}
Read \cref{fig:in:proprietary-market-makers-and-high-frequency-firms:cov}: which fields does the listed firm fill, and
which one does no private firm fill?
\end{exercise}

\begin{exercise}[$\star\star\star$]\label{exo:in:proprietary-market-makers-and-high-frequency-firms:7}
\emph{Coding.} Add to the facts a UK subsidiary's average number of employees as an entity-scope fact. Show that
\texttt{coverage} does not change, and argue from \cref{ex:in:proprietary-market-makers-and-high-frequency-firms:jse} why
it must not.
\end{exercise}

\begin{exercise}[$\star\star\star$]\label{exo:in:proprietary-market-makers-and-high-frequency-firms:8}
\emph{Find the flaw.} ``The firm's European subsidiary earned \$643 million in 2025; Europe is about a quarter of world
trading, so the firm earned about \$2.6 billion.''
\end{exercise}

\section{Problem: What Can the Public Know?}

\begin{problem}[{Weekend problem --- what can the public know?}]\label{pb:in:proprietary-market-makers-and-high-frequency-firms:1}
A student with offers from two private trading firms and one listed firm wants to compare them from public sources
only. You build the profiles.

\medskip\noindent\textbf{Part I --- The records.}
\begin{enumerate}
\item Name the four kinds of public record about a private trading firm.
\item Which part of a US broker-dealer's annual report is public, and which may be kept confidential?
\item What does a UK subsidiary's strategic report describe, and why is its revenue not the firm's?
\item Define a \omterm{def:in:proprietary-market-makers-and-high-frequency-firms:profile}{firm profile} and its scope.
\item Give the five steps of \cref{met:in:proprietary-market-makers-and-high-frequency-firms:build}.
\end{enumerate}

\medskip\noindent\textbf{Part II --- The numbers that exist.}
\begin{enumerate}[resume]
\item Give the subsidiary group's revenue in 2024 and 2025 and its change.
\item Give its after-tax margin in each year and its return on average equity in 2025.
\item Give the broker-dealer's member's capital and its assets-to-capital ratio at the end of 2023.
\item Give the listed firm's 2025 revenue, headcount and revenue per head.
\item Give the share of the exchange group's 2025 volume that was electronic.
\end{enumerate}

\medskip\noindent\textbf{Part III --- The profiles.}
\begin{enumerate}[resume]
\item Which firms state a founding year, and which years?
\item Write each stated headcount as a range.
\item Why can offices not be compared across the table as printed?
\item How many of the five fields does each firm fill at the firm level?
\item Which field does no private firm fill, and why?
\end{enumerate}

\medskip\noindent\textbf{Part IV --- The verdict.}
\begin{enumerate}[resume]
\item State the \emph{named result}: the fields filled by the listed firm, the median for the private firms, and the
number of private firms with a firm revenue.
\item In the race study, what share of races do the top three and the top six firms win?
\item What does the concentration of race winners imply for a specialist's job market?
\item Which two public records would most improve the private firms' profiles, and which of them exists?
\item In two sentences, what should the student conclude about comparing the offers on public numbers?
\end{enumerate}
\end{problem}

\section{Interview questions}

\begin{interviewq}[$\star$ \iqroles{trader, developer}]\label{iq:in:proprietary-market-makers-and-high-frequency-firms:1}
What does it mean for a firm to be a \omterm{def:in:a-map-of-the-industry:ptf}{principal trading firm}, and how does that change what a trader there is paid
for?
\end{interviewq}

\begin{interviewq}[$\star$ \iqroles{researcher}]\label{iq:in:proprietary-market-makers-and-high-frequency-firms:2}
A firm's site says it has ``more than 3\,000 employees''. How would you record that in a dataset?
\end{interviewq}

\begin{interviewq}[$\star\star$ \iqroles{risk}]\label{iq:in:proprietary-market-makers-and-high-frequency-firms:3}
A broker-dealer has assets of 11 times its capital. Is that high for a market maker? What would you want to know first?
\end{interviewq}

\begin{interviewq}[$\star\star$ \iqroles{trader, researcher}]\label{iq:in:proprietary-market-makers-and-high-frequency-firms:4}
Six firms win most latency races and also lose most of them. What does that tell you about the business?
\end{interviewq}

\begin{interviewq}[$\star\star$ \iqroles{developer}]\label{iq:in:proprietary-market-makers-and-high-frequency-firms:5}
Why might a trading firm choose a less common programming language as its primary platform, and what does it cost?
\end{interviewq}

\begin{interviewq}[$\star\star\star$ \iqroles{researcher}]\label{iq:in:proprietary-market-makers-and-high-frequency-firms:6}
A subsidiary's revenue fell 35\% in a year in which the markets it trades were busier than the year before. Give three
explanations that do not involve the firm doing worse.
\end{interviewq}
